% ==============================================================
% CAPÍTULO 2 — METODOLOGIA
% ==============================================================

\chapter{Metodologia}
\label{cap:metodologia}

Este capítulo descreve os procedimentos de coleta de dados e os frameworks analíticos empregados na construção do diagnóstico e da proposta apresentados nos capítulos seguintes.

\section{Coleta de Dados e Pesquisa de Campo}

A base de informações do estudo foi construída a partir da combinação de dados primários e secundários, permitindo o embasamento empírico e teórico das análises.

\subsection{Visita técnica e observação direta}

O grupo realizou uma visita presencial à sede da Capim, em São Paulo. Essa imersão permitiu a observação direta das instalações e o mapeamento preliminar do fluxo de trabalho das equipes e de como elas se dividem no dia a dia.

\subsection{Entrevista estruturada}

Conduziu-se uma entrevista estruturada, por videoconferência, com Lucas Mathias, integrante do time fundador da Capim e responsável pela frente de experimentos e produtos financeiros mais recentes da empresa \cite{capimentrevista2026}. O roteiro contemplou questões sobre as ferramentas de trabalho utilizadas, a trajetória competitiva da empresa, os serviços em constante aprimoramento, os mecanismos de avaliação e continuidade de projetos, a gestão da incerteza e as maiores oportunidades e ameaças percebidas. A transcrição integral da entrevista está disponível no repositório do trabalho e um roteiro-síntese organizado por tema consta no Apêndice~\ref{apend:entrevista}. Esse contato primário foi essencial para a compreensão do modelo de negócio, da cultura interna e, especialmente, da estrutura real do processo decisório de inovação -- objeto central deste relatório.

\subsection{Levantamento de dados secundários}

Realizou-se um levantamento sobre o mercado odontológico brasileiro, os principais concorrentes da Capim (nacionais) e as tecnologias vigentes no setor, a partir de relatórios setoriais do Distrito, da ABStartups, do Conselho Federal de Odontologia (CFO), da Grand View Research, da Agência Nacional de Saúde Suplementar (ANS) e da pesquisa conjunta ABCD/PwC sobre crédito por fintechs, a fim de contextualizar o posicionamento da empresa no setor (Capítulo~\ref{cap:meio-ambiente}).

\section{Frameworks de Análise}

Com os dados empíricos estruturados, aplicou-se um conjunto de frameworks da engenharia de produção e da administração para caracterizar o modelo organizacional da Capim e, em particular, para diagnosticar seu processo decisório de inovação -- foco analítico deste relatório, conforme orientação recebida na devolutiva da versão anterior do trabalho.

\subsection{Frameworks de caracterização organizacional}

\begin{description}
  \item[Estrela de Galbraith] utilizada para examinar o alinhamento entre cinco dimensões da empresa -- estratégia, estrutura, processos, recompensas e pessoas --, buscando validar se a infraestrutura interna sustenta a proposta de valor da Capim (Seção~\ref{sec:galbraith}).
  \item[Configuração de Mintzberg] aplicada para classificar a tipologia estrutural da organização, identificando o mecanismo de coordenação predominante e a parte-chave da organização (Seção~\ref{sec:mintzberg}).
  \item[Análise SWOT] desenvolvida para sintetizar o ambiente competitivo, cruzando forças e fraquezas internas com oportunidades e ameaças externas (Seção~\ref{sec:swot}).
\end{description}

\subsection{Frameworks de processo decisório de inovação}

Para responder diretamente à orientação de aprofundar a análise do sistema decisório com os métodos da disciplina, o Capítulo~\ref{cap:processo-decisorio} aplica quatro frameworks complementares:

\begin{description}
  \item[Tipologia de processos de inovação] de \textcite{salerno2015}, baseada em teoria contingencial, utilizada para classificar cada linha de produto da Capim segundo o tipo de processo de inovação que efetivamente governa seu desenvolvimento, em vez de assumir \textit{a priori} um único modelo linear de idealização ao lançamento.
  \item[Cadeia de valor da inovação] de \textcite{hansenbirkinshaw2007}, empregada para diagnosticar em qual dos três elos -- geração de ideias, conversão ou difusão -- está o principal gargalo do sistema de inovação da Capim.
  \item[Gestão de portfólio] de \textcite{goffinmitchell2010}, cuja matriz de risco e retorno é utilizada para mapear o balanceamento entre projetos incrementais e radicais na carteira atual da empresa.
  \item[Ambidestria e modelo Descoberta--Incubação--Aceleração] discutidos por \textcite{salernogomes2018}, utilizados para avaliar se a Capim trata inovação incremental e radical com sistemas de gestão diferenciados, e para fundamentar a recomendação de um \textit{learning plan} formal \cite{riceoconnorpierantozzi2008} nas etapas de maior incerteza.
\end{description}

A combinação desses frameworks permite ir além da descrição do processo decisório relatado pelo entrevistado -- apresentando-o, no Capítulo~\ref{cap:diagnostico-proposta}, à luz da literatura, com identificação explícita de pontos negativos e de uma proposta de sistema aprimorado.
